\chapter{Conclusion generale}

Ce memoire a repondu a la problematique initiale : identifier des opportunites
de trading robustes dans un marche bruite tout en reduisant les biais de
sur-ajustement. La reponse proposee n'est pas un nouvel indicateur isole, mais
une architecture de decision quantitative en quatre stations.

Les contributions principales sont :
\begin{enumerate}
\item une structuration methodologique Data $\rightarrow$ Signaux
$\rightarrow$ Strategie $\rightarrow$ Backtest ;
\item des contrats inter-stations explicites pour tracer les hypotheses et les
sorties ;
\item une validation orientee robustesse (WFO, OOS, criteres statistiques) ;
\item un socle logiciel extensible pour l'industrialisation progressive.
\end{enumerate}

Au plan metier, la plateforme apporte une aide a la decision plus rigoureuse
qu'un backtest monolithique traditionnel. Au plan scientifique, elle impose une
discipline de separation entre recherche, parametrage et validation, conforme aux
recommandations de la litterature.

Les prochaines etapes consistent a elargir l'univers de validation, renforcer la
detection de regime et formaliser un protocole de passage en pre-production.

